\documentclass[11pt]{article}
\usepackage[margin=1in]{geometry}
\usepackage{natbib}
\usepackage{booktabs}
\usepackage{amsmath}
\usepackage{graphicx}
\usepackage{hyperref}
\usepackage{lineno}

\title{Stable population structure in Europe since the Iron Age, despite high mobility}
\author{
Antonio, M.L.$^{1,\dagger}$ \and Weiß, C.L.$^{1,\dagger}$ \and Gao, Z.$^{2,\dagger}$ \and Moots, H.M.$^{1}$ \and Rickards, O.$^{3}$ \and Mathieson, I.$^{1,*}$\\[6pt]
\small $^{1}$Department of Genetics, Perelman School of Medicine, University of Pennsylvania, Philadelphia, PA, USA\\
\small $^{2}$Department of Biology, University of Pennsylvania, Philadelphia, PA, USA\\
\small $^{3}$Centre of Molecular Anthropology for the Study of Ancient DNA, Department of Biology, University of Rome ``Tor Vergata'', Rome, Italy\\
\small $^{\dagger}$These authors contributed equally\\
\small $^{*}$Corresponding author: \texttt{mathi@pennmedicine.upenn.edu}
}
\date{}

\begin{document}
\maketitle

\begin{abstract}
Ancient DNA research in the past decade has revealed that European population structure changed dramatically in the prehistoric period (14,000--3,000 years before present, YBP), reflecting the widespread introduction of Neolithic farmer and Bronze Age Steppe ancestries. However, little is known about how population structure changed from the historical period onward (3,000 YBP -- present). To address this, we collected whole genomes from 204 individuals from Europe and the Mediterranean, many of which are the first historical period genomes from their region (e.g.\ Armenia and France). We found that most regions show remarkable inter-individual heterogeneity. At least 7\% of historical individuals carry ancestry uncommon in the region where they were sampled, some indicating cross-Mediterranean contacts. Despite this high level of mobility, overall population structure across western Eurasia is relatively stable through the historical period up to the present, mirroring geography. We show that, under standard population genetics models with local panmixia, the observed level of dispersal would lead to a collapse of population structure. Persistent population structure thus suggests a lower effective migration rate than indicated by the observed dispersal. We hypothesize that this phenomenon can be explained by extensive transient dispersal arising from drastically improved transportation networks and the Roman Empire's mobilization of people for trade, labor, and military. This work highlights the utility of ancient DNA in elucidating finer-scale human population dynamics in recent history.
\end{abstract}

\section{Introduction}

Ancient DNA (aDNA) sequencing has provided immense insight into previously unanswered questions about human population history \citep{haak2015, lazaridis2014, fu2016, allentoft2015}. Initially, sequencing efforts focused on identifying the main ancestry groups and transitions during prehistoric times, for which there is no written record. Key discoveries include the identification of three major ancestral populations contributing to present-day Europeans: Western Hunter-Gatherers, Early Neolithic Farmers, and Bronze Age Steppe pastoralists \citep{lazaridis2014, haak2015}. Massive migrations during the Neolithic and Bronze Age fundamentally reshaped European genetic structure \citep{allentoft2015, brace2019, cassidy2016}.

Recently, aDNA sampling has expanded to more recent historical times, allowing the study of movements of people using genetic data alongside the well-studied historical record \citep{antonio2019, olalde2018, ebenesersdottir2018, margaryan2020}. Ancient Rome, for example, was shown to be a genetic crossroads of Europe and the Mediterranean \citep{antonio2019}. However, we lack a comprehensive assessment of historical genetic structure, including characterizing genetic heterogeneity and interactions across regions. Integrating historical period genetics is instrumental to better understanding the development of European and Mediterranean population structure from prehistoric to present-day.

Prehistoric ancient DNA studies have documented dramatic ancestry turnovers, including the spread of Neolithic farmer ancestry from Anatolia \citep{broushaki2016, brace2019} and the later introduction of Steppe ancestry \citep{haak2015, damgaard2018, debarros2018}. The genetic landscape of Europe was fundamentally transformed by these events. Yet, relatively little is known about the genetic consequences of the well-documented historical migrations and movements during the Iron Age, Roman period, and Late Antiquity, despite the abundance of historical and archaeological evidence for substantial population mobility \citep{dejong2016, cherry2007}.

In this study, we address this gap by generating and analyzing 204 new ancient genomes from across Europe and the Mediterranean, spanning the historical period (approximately 1000 BCE to 1000 CE). We integrate these with published prehistoric and historical genomes to characterize regional population structure and heterogeneity, identify ancestry outliers indicative of long-range mobility, and examine the temporal stability of spatial genetic structure. Through simulation, we demonstrate that the observed levels of long-range dispersal should, under standard models, erode population structure far more than what is observed, pointing to a distinction between observed and effective migration rates.

\section{Results}

\subsection{204 new historical genomes from Europe and the Mediterranean}

We collected whole genomes from 204 individuals across 53 archaeological sites in 18 countries spanning Europe and the Mediterranean, 26 of which were recently reported \citep{moots2022}. Notably, these include the first historical period genomes (after 1000 BCE) from present-day Armenia, Algeria, Austria, and France. Dates for 126 samples were directly determined through radiocarbon dating, and were used alongside archaeological contexts to infer dates for the remaining samples.

DNA was extracted from either the powdered cochlear portion of the petrous bone ($n = 203$) or from teeth ($n = 1$). We performed whole genome sequencing to a median depth of 0.92$\times$ (range: 0.16$\times$ to 2.38$\times$). For downstream integration with published data, pseudohaploid genotypes were called for the 1240k SNP panel \citep{mathieson2015}, resulting in a median of 685,058 SNPs (range: 167,000 to 1,029,345) per sample.

We analyzed newly reported genomes in conjunction with 2,033 present-day genomes, 1,998 prehistoric genomes, and 764 published historical period genomes from multiple sources \citep{clemente2021, kovacevic2014, mallick2023, pagani2016, saupe2021, zegarac2021}. Genomes were grouped by regions and time periods and analyzed using principal component analysis (PCA) and qpAdm modeling \citep{haak2015, patterson2012}.

\subsection{Local historical population structure varies across regions}

To investigate historical population structure, we categorized the data into spatio-temporal groups and characterized inter-individual heterogeneity within these groups by examining: (1) variation of projections onto a PCA space of present-day genomes, (2) genetic groups identified by qpAdm and clustering across time within a region, and (3) admixture profiles of those groups.

PCA was performed on 829 individuals (480,712 SNPs) using smartpca v1600. A majority of regions have highly heterogeneous populations in at least some historical time periods. This heterogeneity is illustrated by both the visual spread in PCA and the genetically distinct clusters of individuals based on pairwise modeling with qpAdm \citep{haak2015, harney2021}. On average, we identified 10 genetic clusters within each region present during the historical period, with a minimum of 2 and a maximum of 23.

In Armenia, for example, the population is highly homogeneous at any given time. After the Copper Age, there are two distinct genetic clusters, separated by a temporal split around 772--403 BCE. The earlier cluster (C1) includes newly reported samples ($n = 5$) from Beniamin and published ones ($n = 6$) from five other sites. The later cluster (C3), which contains newly reported samples ($n = 12$) from Beniamin dating between 403 BCE--500 CE, is genetically similar to present-day Armenians. Despite the split, there is evidence of partial continuity: the later cluster (C3) can be modeled using around 50\% of the earlier cluster (C1) and an additional source of Steppe ancestry.

Historical genomes from Northern Europe, particularly newly reported genomes from Lithuania and Poland, exhibit a similar level of homogeneity. In contrast, newly collected samples reinforce previous findings of high heterogeneity in Rome, including a large portion of the population having affinities for present-day Near Eastern populations \citep{antonio2019, posth2021}.

Interestingly, Southeastern European and Western European individuals during the Imperial Roman and Late Antiquity period also exhibit high heterogeneity, on par with that of contemporaneous Italy. In Southeastern Europe, a core group of individuals have ancestry similar to that of present-day and contemporaneous Central Europeans, while other clusters have ancestry similar to that of Northern Europeans and Eastern Mediterraneans. These ancestry groups are found in contemporaneous Italy and Western Europe as well. We also observe individuals of eastern nomadic ancestry, similar to that of Sarmatian individuals previously reported, in both Western Europe ($n = 2$) and Southeastern Europe ($n = 2$).

\subsection{Ancestry outliers and long-range mobility}

Overall, we see remarkable local genetic heterogeneity as well as cross-regional similarities which point to common patterns of mobility. We considered an individual an outlier if they belonged to an ancestry cluster that is underrepresented (consisting of fewer than 5\% of individuals in a region or at most two individuals) within their sampling region from the Bronze Age up to present-day. To focus on first-generation migrants as well as long-range movements, we further identified outlier individuals who can be modeled as 100\% (i.e., a one-component model) of a majority ancestry cluster found in a different region.

In total, we identified 11\% of individuals as outliers, and could connect 7\% of individuals to a putative source in a different region. Based on the regions where these outliers and their sources originated, we created a network to illustrate their movements. For example, the Armenian population is quite homogeneous, and unsurprisingly, no outliers were found within Armenia. However, we found outlier individuals in the Levant and Italy who can be putatively traced back to Armenia.

In North Africa, outliers found in Iron Age Tunisia \citep{moots2022} indicate movements from many regions in Europe, and North African-like outliers were found in Italy and Austria (Western Europe). North African ancestry in Italy is supported by a single previously reported individual from the Imperial Roman period \citep{antonio2019}. Similar North African ancestry in Western Europe is supported by a single individual from Wels, Austria, a site located on the frontier of the Roman Empire. This individual from Austria can be modeled using Canary Islander individuals from the Medieval Ages or an Iron Age outlier from Kerkouane, a Punic city near Carthage in modern-day Tunisia.

The 7\% estimate for outliers with source should be considered a lower bound. There are several cases where a cluster comprises more than 5\% of the individuals in the region, but are clearly of a different ancestry than the majority and seem to be transient. We expect the actual proportion of individuals involved in long-distance movements to be higher than reported here.

\subsection{Spatial population structure is relatively stable in the last 3,000 years}

The remarkable amount of heterogeneity and mobility in the historical period leads to the question of what effect this had on overall spatial population structure. We computed pairwise $F_{ST}$ \citep{bhatia2013} between regional groups in each time period. In each time period, we recovered the classical pattern of isolation-by-distance, where individuals closer in geographic space are also more similar genetically.

Across time periods, we see a large decrease in overall $F_{ST}$ from the Mesolithic and Neolithic periods to the Bronze Age (approximately 10,000--2300 BCE), coinciding with the major prehistoric migrations \citep{haak2015, lazaridis2014}. From the Bronze Age onward, however, $F_{ST}$ does not decrease further with time, indicating that the level of genetic differentiation across space is relatively stable from the Bronze Age to present-day. Confidence intervals were calculated through a bootstrap procedure, using 200 bootstrap replicates.

To further visualize spatial structure, we performed PCA on 829 present-day European and Mediterranean genomes sampled across geographical space and projected historical period genomes onto the same PC space. Echoing close correspondence between genetic structure and geographic space in present-day Europeans \citep{novembre2008}, we recovered similar spatial structure for historical samples as well, although noisier due to a narrower sampling distribution and higher local genetic heterogeneity. The similarity in structure between present-day and the historical period is especially striking in comparison to a projection of prehistoric genomes, which shows much weaker correspondence to the present-day PCA as well as to geographic space.

Together, our analyses indicate that European and Mediterranean population structure has been relatively stable over the last 3,000 years.

\subsection{Simulations of population structure under long-range dispersal}

This raises the question: is it surprising for stable population structure to be maintained in the presence of approximately 7--11\% long-range migration? To address this, we performed spatial simulations using a stepping-stone model.

In these simulations, spatial population structure is established through local mate choice and limited dispersal, which we calibrated to approximately match the spatial differentiation observed in historical-period Europe (maximum $F_{ST}$ of approximately 0.03). We used the parameter pair $N = 50{,}000$ and $\sigma_{\text{Disp}} = 0.02$ for all simulations. We then allowed a proportion of the population to disperse longer distances, empirically matching the migration distances observed in the data during the historical period.

Even with long-range dispersal as low as 4\%, we observe decreasing $F_{ST}$ over 120 generations (approximately 3,000 years with a generation time of 25 years) as individuals become less differentiated genetically across space. At 8\%, $F_{ST}$ decreases dramatically within 120 generations as spatial structure collapses to the point that it is hardly detectable in the first two principal components.

These simulations indicate that under a basic spatial population genetics model we would expect structure to collapse by present-day given the levels of movement we observe. Persistent population structure thus suggests a lower effective migration rate than indicated by the observed dispersal, which we hypothesize is explained by extensive transient dispersal---individuals who moved long distances but did not contribute proportionally to the local gene pool.

\section{Discussion}

Our analysis of 204 new ancient genomes from across Europe and the Mediterranean, integrated with nearly 5,000 published genomes, reveals a paradox at the heart of historical population genetics: despite remarkable levels of individual mobility during the Iron Age, Roman period, and Late Antiquity, the overall spatial structure of European populations has remained relatively stable for the past 3,000 years.

The high levels of inter-individual heterogeneity and the identification of at least 7--11\% of individuals as ancestry outliers demonstrate that long-range movement was common during the historical period. This is consistent with historical evidence for the extensive mobilization of people throughout the Roman Empire for trade, military service, and forced labor \citep{dejong2016, cherry2007, erdkamp2015, wardperkins2007}. The network of outlier movements we reconstruct mirrors known historical patterns of connectivity across the Roman world.

The stability of population structure in the face of this mobility is initially surprising. Our simulations demonstrate that, under standard population genetics models with local panmixia, even modest rates of long-range dispersal (4--8\%) would lead to a dramatic collapse of spatial structure over 120 generations. The fact that structure has persisted suggests that the effective migration rate---the rate at which migrants contribute to the local gene pool---is substantially lower than the observed rate of dispersal.

We hypothesize that this discrepancy can be explained by transient dispersal: many individuals who moved long distances did not permanently settle and reproduce in their destination populations. The Roman Empire's transportation networks facilitated unprecedented mobility, but much of this movement may have been temporary, involving soldiers, traders, and laborers who eventually returned home or did not integrate reproductively into local populations. This interpretation is consistent with archaeological and historical evidence for transient populations in Roman frontier towns and cities.

Our findings demonstrate the power of ancient DNA for revealing fine-scale population dynamics in recent historical periods. While the broad strokes of European genetic history were painted during the Neolithic and Bronze Age transitions, the historical period added considerable local complexity through mobility without fundamentally altering the continental-scale structure. This work also highlights the importance of dense sampling across both space and time to capture the full range of human genetic diversity in the past, and the need for population genetics models that account for the distinction between physical mobility and effective gene flow.

\section{Materials and Methods}

\subsection{Sample collection and archaeological sites}

Sampling was performed to maximize coverage across Europe and the Mediterranean, with particular focus on regions where no published genomic data existed for the Imperial Roman and Late Antiquity periods at the time of sampling (e.g., France, Austria, the Balkans, Armenia, and North Africa). Although we aimed to collect samples in the Imperial Roman and Late Antiquity period (approximately 1 CE--700 CE), some samples fall outside of this period due to limited sample availability.

\subsection{Date determination and time periods}

Determination of time periods was influenced by: (1) the wide geographic range represented in the data, (2) the amount of data available in the proposed time periods, (3) the amount of genetic change observed during a time interval, and (4) the types of temporal comparisons made in the analysis. The following time periods were used: Mesolithic and Neolithic (10000--3500 BCE), Copper Age (3500--2300 BCE), Bronze Age (2300--1000 BCE), Iron Age (1000 BCE--1 CE), Imperial Roman and Late Antiquity (1--700 CE), and Medieval to Present (700 CE--present).

\subsection{DNA extraction and sequencing}

DNA was extracted from the powdered cochlear portion of the petrous bone ($n = 203$) or from teeth ($n = 1$) using established protocols \citep{dabney2013, rohland2015}. Whole genome sequencing was performed to a median depth of 0.92$\times$. Pseudohaploid genotypes were called for the 1240k SNP panel \citep{mathieson2015}.

\subsection{Population genetic analyses}

PCA was performed on 829 present-day individuals (480,712 SNPs) using smartpca v1600, with 5 outlier iterations, 10 principal components along which to remove outliers, altnormstyle set to NO, and least-squares projection enabled. Ancient genomes were projected onto the present-day PCA space.

Ancestry modeling was performed using qpAdm \citep{haak2015, patterson2012}. Pairwise $F_{ST}$ was computed following the method of \citet{bhatia2013} with 200 bootstrap replicates for confidence intervals.

\subsection{Simulation of spatial population structure}

We simulated spatial population structure using a stepping-stone model calibrated to match the observed maximum $F_{ST}$ of approximately 0.03. A grid search across parameter pairs yielded $N = 50{,}000$ and $\sigma_{\text{Disp}} = 0.02$. Long-range dispersal was introduced at rates of 4\% and 8\%, with dispersal distances calibrated to match the distribution of outlier distances observed in the data ($\sigma_{\text{DispLR}} = 0.20$). The effect of long-range dispersal on spatial $F_{ST}$ was tracked over 120 generations.

\bibliographystyle{plainnat}
\bibliography{references}

\end{document}
